\documentclass[11pt]{article}
\usepackage[a4paper,margin=1in]{geometry}
\usepackage{amsmath,amssymb}
\usepackage{booktabs}
\usepackage{array}
\usepackage{xcolor}

\newcommand{\Z}{\mathbb{Z}}
\newcommand{\md}{\bmod\ 2^{32}}
\newcommand{\code}[1]{\texttt{#1}}

% Lightweight dependency-free pseudocode block.
\newenvironment{pseudo}
  {\begin{quote}\ttfamily\begin{tabbing}
   \hspace{2em}\=\hspace{2em}\=\hspace{2em}\=\kill}
  {\end{tabbing}\end{quote}}

\title{\vspace{-1.5em}The Tree--Walk Hashing Kernel\\[0.3em]
       \large Reference algorithm and optimized implementation}
\author{}
\date{}

\begin{document}
\maketitle
\vspace{-3em}

\noindent\textbf{Part I} states the algorithm exactly as the grader's
reference implements it. \textbf{Part II} explains how the optimized kernel
runs the same computation in \textbf{1{,}053} simulated cycles instead of the
shipped baseline's \textbf{147{,}734} (a $140\times$ reduction), broken into
three independent concerns: \emph{algorithmic} changes (what is computed),
\emph{slot allocation} (how operations are packed onto the machine's
parallel engines each cycle), and \emph{memory allocation} (how the
1{,}536-word scratch space is used).

\begin{center}\rule{0.4\linewidth}{0.4pt}\end{center}
\section*{\centering Part I --- The reference (naive) algorithm}

\section{Overview}

The workload is a fixed-length, data-dependent walk over a static binary
tree, run independently for a batch of walkers. Each walker repeatedly reads
the value at its current tree node, mixes it into its running state with a
32-bit hash, and uses one bit of the new state to choose which child to
descend to. All arithmetic is over the ring $\Z_{2^{32}}$: every operation is
reduced modulo $2^{32}$.

\section{Data}

\paragraph{Tree.} An implicit perfect binary tree of height $H$, stored as a
flat array $T[0\,\dots\,N-1]$ of $N = 2^{\,H+1}-1$ node values. Node $x$ has
children $2x+1$ (left) and $2x+2$ (right); the root is index $0$. Each
$T[x]$ is a fixed value in $[0,\,2^{30}-1]$.

\paragraph{Walkers.} A batch of $B$ walkers, each with an index
$I[i]\in\{0,\dots,N-1\}$ and a state value $V[i]\in\Z_{2^{32}}$. Initially
every walker sits at the root, $I[i]=0$, and $V[i]$ is a fixed value in
$[0,\,2^{30}-1]$. The walk runs for $R$ rounds.

\section{The hash function}

$\mathrm{myhash}:\Z_{2^{32}}\to\Z_{2^{32}}$ is applied in six stages. Writing
$a$ for the value entering a stage, each stage recomputes
\[
  a \;\leftarrow\; \bigl((a \mathbin{\square_1} v_1)\ \square_2\ (a \mathbin{\square_3} v_3)\bigr)\ \md,
\]
where the two half-expressions are each reduced modulo $2^{32}$ before being
combined by $\square_2$. The operators $\square$ are add ($+$), xor
($\oplus$), left shift ($\ll$), or right shift ($\gg$). The six stages use:

\begin{center}
\begin{tabular}{cccccc}
\toprule
stage & $\square_1$ & $v_1$ & $\square_2$ & $\square_3$ & $v_3$ \\
\midrule
1 & $+$      & \code{0x7ED55D16} & $+$      & $\ll$ & $12$ \\
2 & $\oplus$ & \code{0xC761C23C} & $\oplus$ & $\gg$ & $19$ \\
3 & $+$      & \code{0x165667B1} & $+$      & $\ll$ & $5$  \\
4 & $+$      & \code{0xD3A2646C} & $\oplus$ & $\ll$ & $9$  \\
5 & $+$      & \code{0xFD7046C5} & $+$      & $\ll$ & $3$  \\
6 & $\oplus$ & \code{0xB55A4F09} & $\oplus$ & $\gg$ & $16$ \\
\bottomrule
\end{tabular}
\end{center}

Explicitly, with $C_k$ the constant of stage $k$ (both halves read the same
input $a$; the update is not sequential within a stage):
\begin{align*}
  a &\leftarrow (a + C_1) + (a \ll 12), &
  a &\leftarrow (a \oplus C_2) \oplus (a \gg 19), &
  a &\leftarrow (a + C_3) + (a \ll 5), \\
  a &\leftarrow (a + C_4) \oplus (a \ll 9), &
  a &\leftarrow (a + C_5) + (a \ll 3), &
  a &\leftarrow (a \oplus C_6) \oplus (a \gg 16).
\end{align*}

\section{The per-round update and full algorithm}

For a single walker with index $x$, value $v$, one round reads $n=T[x]$, sets
$v \leftarrow \mathrm{myhash}(v \oplus n)$, descends
$x \leftarrow 2x + 1 + (v \bmod 2)$ (left child if $v$ even, right if odd),
and wraps $x \leftarrow 0$ if $x \ge N$. The reference iterates rounds outside,
walkers inside:

\begin{pseudo}
\textbf{for} $h \leftarrow 1$ \textbf{to} $R$: \\
\> \textbf{for} $i \leftarrow 1$ \textbf{to} $B$: \\
\> \> $n \leftarrow T[\,I[i]\,]$ \\
\> \> $V[i] \leftarrow \mathrm{myhash}\bigl(V[i] \oplus n\bigr)$ \\
\> \> $I[i] \leftarrow 2\,I[i] + 1 + (V[i] \bmod 2)$ \\
\> \> \textbf{if } $I[i] \ge N$ \textbf{ then } $I[i] \leftarrow 0$ \\
\end{pseudo}

\noindent The graded output is the final value array $V[0\,\dots\,B-1]$; the
indices are internal bookkeeping. The benchmark instance is $H=10$
($N=2047$), $B=256$, $R=16$ --- $B\cdot R = 4096$ hash evaluations. The tree
has 11 levels; a walker at level $\ell$ moves to level $\ell+1$ each round
(wrapping from the bottom), so over 16 rounds it visits levels
$0,1,\dots,10,0,1,2,3,4$.

\begin{center}\rule{0.4\linewidth}{0.4pt}\end{center}
\section*{\centering Part II --- The optimized implementation}

\section{The target machine}

The kernel is scored on a single-core VLIW/SIMD simulator. Each cycle issues
one \emph{bundle} of \emph{slots} spread across six engines, up to a
per-engine limit; the cycle count is the number of non-trivial bundles.
Vector ops act on \code{VLEN}${}=8$ lanes at once. Scratch is a flat
1{,}536-word array that serves as \emph{both} the register file and working
memory --- there is no separate register space; \code{mem} is a large
separate array reached only through the load/store engines.

\begin{center}
\begin{tabular}{lcp{0.60\linewidth}}
\toprule
engine & slots/cyc & what it does \\
\midrule
\code{alu}   & 12 & scalar integer ops (one lane each) \\
\code{valu}  & 6  & vector ops, 8 lanes/slot; \textbf{the only engine with fused \code{multiply\_add}} ($a\!\cdot\!b+c$) \\
\code{load}  & 2  & \code{mem} read: scalar gather \code{mem[scratch[addr]]}, or \code{vload} of 8 contiguous words \\
\code{store} & 2  & \code{mem} write (scalar or 8-wide \code{vstore}) \\
\code{flow}  & 1  & \code{select}/\code{vselect} (8-lane mux) and \code{add\_imm} only --- no arithmetic \\
\code{debug} & 64 & correctness compares; free (never counted in cycles) \\
\bottomrule
\end{tabular}
\end{center}

\noindent We measure work in two units: a \emph{slot} is one engine
issue-slot in one cycle; a \emph{lane-op} is one scalar-equivalent operation
(a \code{valu} slot does 8 lane-ops, an \code{alu} slot does 1). The peak
arithmetic throughput is $12 + 6\!\cdot\!8 = 60$ lane-ops/cycle across
\code{alu}${}+{}$\code{valu}.

\section{Result and where the cycles go}

\begin{center}
\begin{tabular}{rl}
\toprule
cycles & milestone \\
\midrule
147{,}734 & shipped scalar baseline (1 op/bundle) \\
12{,}414  & vectorize: 8 walkers per \code{valu} op \\
4{,}990   & software-pipeline the phases \\
2{,}148   & fused hash + carry gather address + gather-under-compute \\
1{,}572   & tournament routing of tree levels 1--3 (fewer gathers) \\
1{,}140   & group skewing (overlap compute- and load-heavy blocks) \\
1{,}053   & \textbf{current} (constant commuting, engine racing, drain/ramp fixes) \\
\bottomrule
\end{tabular}
\end{center}

At 1{,}053 cycles the two arithmetic engines are jointly saturated and the
schedule is within $1.7\%$ of the binding resource floor:

\begin{center}
\begin{tabular}{lrrr}
\toprule
engine & slots used & capacity & utilization \\
\midrule
\code{valu}  & 6{,}209  & 6{,}318  & 98.3\% (floor $6209/6 = 1035$ cyc) \\
\code{alu}   & 12{,}169 & 12{,}636 & 96.3\% \\
\code{load}  & 1{,}924  & 2{,}106  & 91.4\% (1{,}875 gathers) \\
\code{flow}  & 637      & 1{,}053  & 60.5\% \\
\code{store} & 38       & 2{,}106  & 1.8\% \\
\bottomrule
\end{tabular}
\end{center}

\noindent Splitting that same slot count by \emph{purpose} instead of just
by engine shows where each engine's cycles actually go, and lines it up
against the \emph{lane-op} total (the scalar-equivalent work, after
un-packing \code{valu}/\code{flow}'s 8-wide slots) so the two units --- one
slot costs one cycle-slot on its engine, one lane-op is one scalar
operation --- can be read side by side:

\begin{center}
\begin{tabular}{lrrrrrrr}
\toprule
purpose & \code{alu} & \code{valu} & \code{load} & \code{store} & \code{flow} & lane-ops & \% \\
\midrule
Hash    & 10{,}280 & 4{,}547 & ---     & ---   & ---   & 46{,}656 & 67.8\% \\
Routing &    809   &   593   & 1{,}872 & ---   & 623   & 12{,}409 & 18.0\% \\
Index   &  1{,}064 & 1{,}010 & ---     & ---   & ---   &  9{,}144 & 13.3\% \\
Setup   &     16   &    59   &    52   & ---   &  14   &    554  &  0.8\% \\
Store   &    ---   &   ---   &   ---   &  38   & ---   &     38  &  0.1\% \\
\midrule
total   & 12{,}169 & 6{,}209 & 1{,}924 &  38   & 637   & 68{,}801 & 100.0\% \\
\bottomrule
\end{tabular}
\end{center}

\noindent Each column foots to the engine-utilization table above:
\code{alu} and \code{valu} are almost entirely Hash; \code{load} and
\code{flow} are almost entirely Routing (the gather-vs-tournament trade);
\code{store} is 100\% end-of-run writes.

The binding engine is \code{valu} (its fused multiply-add has no scalar
substitute), so the optimizations below are ultimately judged by how many
\code{valu} slots they remove or how well they keep all six engines busy in
the same cycle.

\noindent Zooming in one level further, from slots to individual
\emph{opcodes}, shows which instructions each engine's budget is actually
spent on (\code{\^{}}~xor, \code{\&}~and, \code{|}~or, \code{<<}/\code{>>}
shift left/right):

\begin{center}
\begin{tabular}{llrr}
\toprule
engine & opcode & count & \% of engine \\
\midrule
\code{alu}
  & \code{\^{}}          & 6{,}154 & 50.6\% \\
  & \code{>>}            & 4{,}160 & 34.2\% \\
  & \code{\&}            &   720   &  5.9\% \\
  & \code{<<}            &   475   &  3.9\% \\
  & \code{+}             &   410   &  3.4\% \\
  & \code{-}             &   249   &  2.0\% \\
  & \code{|}             &     1   &  0.0\% \\
\midrule
\code{valu}
  & \code{multiply\_add} & 2{,}846 & 45.8\% \\
  & \code{v\^{}}         & 1{,}991 & 32.1\% \\
  & \code{v\&}           &   552   &  8.9\% \\
  & \code{v>>}           &   508   &  8.2\% \\
  & \code{v+}            &   160   &  2.6\% \\
  & \code{v-}            &    93   &  1.5\% \\
  & \code{vbroadcast}    &    59   &  0.9\% \\
\midrule
\code{load}
  & \code{load}          & 1{,}875 & 97.5\% \\
  & \code{vload}         &    40   &  2.1\% \\
  & \code{const}         &     9   &  0.5\% \\
\midrule
\code{store}
  & \code{vstore}        &    38   & 100.0\% \\
\midrule
\code{flow}
  & \code{vselect}       &   623   & 97.8\% \\
  & \code{add\_imm}      &    12   &  1.9\% \\
  & \code{pause}         &     2   &  0.3\% \\
\midrule
\code{debug} & \code{vcompare} & 384 & 100.0\% (free, uncounted) \\
\bottomrule
\end{tabular}
\end{center}

\noindent Of the ISA's 47 instruction forms across all six engines, the
kernel uses 22. What it uses is narrow and deliberate: every \code{alu}/
\code{valu} op is xor/shift/and-heavy integer mixing (the hash and index
recurrences) plus the fused \code{multiply\_add}; \code{load} is almost
entirely data-dependent gathers; \code{flow} is almost entirely tournament
\code{vselect} folds. \code{load\_offset}, scalar \code{select}, and every
jump/halt/\code{coreid} control-flow op go unused --- the kernel is a
straight-line, branch-free bundle sequence, so it never needs control flow
at all.

\section{Algorithmic improvements}

These change \emph{what} is computed --- the operation count and data-flow
structure --- while leaving the graded output bit-identical.

\paragraph{1. Exploit walker independence (vectorize).} Within a round the
$B=256$ walkers are mutually independent, so they pack 8-to-a-lane onto
\code{valu}: 256 walkers become 32 \emph{groups} of 8. Only the 16 rounds are
sequential. This alone is the $147\text{k}\!\to\!12\text{k}$ step and is the
foundation everything else builds on.

\paragraph{2. Hash fusion.} Three of the six stages ($1,3,5$) have the affine
shape $(a+C)+(a\ll s) = a\,(1+2^s)+C$, which is exactly one
\code{multiply\_add}. Better, stages $2$ and $3$ \emph{compose} across their
boundary: both branches of stage $3$ are affine in stage $2$'s input, so the
pair collapses to two multiply-adds plus one xor rather than four ops. The
6-stage hash reduces from a naive $\sim$18 operations to \textbf{11 mixing
operations} per evaluation. An exhaustive machine search (adjacent-segment
plus meet-in-the-middle, ${\sim}2.4\times10^{12}$ candidates with constants
solved over all $2^{32}$) found no shorter form at any cut point; 11 is very
likely minimal for this composition, though not proven globally.

\paragraph{3. Commute a constant across the round boundary.} Stage 6's
constant xor $\oplus C_6$ can be pushed out of the per-round hash by
pre-transforming the tree once: if the stored node values carry $C_6$
already, the fold-in $v \oplus n$ absorbs the elided $\oplus C_6$ from the
previous round. This removes one operation per evaluation on $9$ of the $16$
rounds. Parity flips by $C_6\bmod 2$ but is corrected for free by reversing
the routing tables and adjusting compile-time offset constants.

\paragraph{4. Tournament routing (the load$\to$arithmetic trade).} The naive
gather $T[I[i]]$ is a data-dependent \code{load}; with one per walker per
round that is $4096/2 = 2048$ load-cycles --- the original wall. But the
\emph{shallow} tree levels are tiny and shared: level $d$ has only $2^d$
distinct nodes, identical for all walkers. Those levels are pre-broadcast
into scratch once, and each walker selects its node with a \emph{tournament}
of $2^d\!-\!1$ multiplexes driven by its accumulated branch-parity bits,
running on \code{valu}/\code{flow} instead of \code{load}. Levels 1--4 are
served this way (a partial fifth level too); levels 5--10 still gather. This
cuts gathers from 4096 to $\sim$1875 and moves the bottleneck off the 2-wide
\code{load} engine onto the arithmetic engines --- the decisive
$2148\!\to\!{\sim}1400$ structural change. A walker's position within a level
is the $d$-bit number formed by its last $d$ parity bits (oldest = most
significant).

\paragraph{5. Position-accumulator index recurrence.} Rather than recompute a
full tree address each round, a walker on tournament levels carries a compact
position $p \leftarrow 2p + b$ (one multiply-add; $b$ is the new parity bit),
and only materializes a real gather address $g = 2p + b + (\text{base} +
2^{L}-1)$ at the boundary where it leaves tournament levels for gather levels.
The ``$+1$'' of $2x{+}1$ and the tree base pointer fold into compile-time
constants. The whole index recurrence costs ${\sim}2.2$ lane-ops per
evaluation --- below the naive 3-op recurrence.

\paragraph{6. Deterministic wraparound.} Because every walker advances exactly
one level per round in lockstep, the ``\,if $x\ge N$ wrap to $0$\,'' branch
fires on the same known rounds for all lanes at once. It is compiled away
entirely --- there is no per-walker compare or select for wraparound; the
position state is simply re-seeded from the root on those rounds.

\section{Slot allocation}

These do not change the computation; they decide \emph{which engine and which
cycle} each operation issues on, to keep the six parallel engines full.

\paragraph{The list scheduler.} Operations are emitted as an unordered pool
and placed greedily at the earliest cycle with (a) all inputs ready and (b) a
free slot on a capable engine, tracking read/write hazards per scratch
address. This replaced a fixed ``wave of six'' hand-schedule and is what lets
the following placement tricks compose.

\paragraph{Multi-encoding engine racing.} Many operations have more than one
legal encoding, and the scheduler places each on whichever engine retires it
soonest:
\begin{itemize}\itemsep2pt
  \item An 8-wide elementwise vector op costs \emph{one} \code{valu} slot or
    \emph{eight} scalar \code{alu} slots. When \code{valu} is backed up and
    \code{alu} has idle lanes, the split wins --- ``\code{alu} offload''
    raising effective throughput toward $60$ lane-ops/cycle.
  \item A tournament fold is either a \code{valu} multiply-add or a
    \code{flow} \code{vselect} (the parity is a raw $0/1$ condition). Each
    fold races the two; ${\sim}\,54\%$ move to the otherwise-idle \code{flow}
    engine, but only where they do not serialize on its single slot.
  \item The index multiply-by-2 recurrence spells either as one \code{valu}
    multiply-add or as \code{shift}${+}$\code{add} across 8 \code{alu} lanes,
    chosen per site.
\end{itemize}
A hard, unconditional version of any of these \emph{loses} (measured): the
\code{flow} engine is 1-wide and its idle time is anti-correlated with when
folds are ready, so blind offloading serializes. Only the per-site race wins.

\paragraph{Cut the operation, not just its placement.} Tournament select
conditions are taken directly from the raw parity vectors (which already
exist for the index update, and ride an otherwise-dead buffer) instead of
being re-extracted with mask-and-shift ops --- eliminating those extractions
from \code{alu}/\code{valu} at zero scratch cost.

\paragraph{Group skewing (software pipeline across blocks).} The 32 groups
are emitted in 4 blocks staggered 3 rounds apart. A round that is
compute-heavy (a tournament level, no gathers) for one block then overlaps a
round that is load-heavy (a deep gather level) for another, so
\code{load} and \code{valu} stay busy in the same cycles instead of taking
turns. This was the single largest scheduling win
($1418\!\to\!1140$).

\paragraph{Setup and drain scheduling.} The startup ramp and final-round
drain are where engines idle. Setup constants are computed on the idle
\code{alu} (e.g. one constant is another shifted) instead of consuming
\code{load} slots; the 32 base-address adds are moved off the 1-wide
\code{flow} engine into parallel \code{alu} chains; disjoint final stores are
paired to use both \code{store} slots; and the last round's tournament fold
is reordered so the last-arriving parity bit selects last, shortening the
drain's critical path.

\section{Memory allocation}

Scratch is only 1{,}536 words and doubles as the register file, so every live
value competes for it. The kernel runs at \textbf{1{,}535/1{,}536} words ---
essentially full --- which is only possible because of the following.

\paragraph{Two tiers.} \code{alu}/\code{valu}/\code{flow} operate directly on
scratch; the large \code{mem} array is reachable only via load/store. Values
that must persist live in scratch; anything that can be recomputed or streamed
from \code{mem} need not.

\paragraph{Persistent state vs.\ reused temp pools.} Each of the 32 groups
keeps exactly three live 8-word vectors across the whole run: its hash state,
its node-value buffer, and its position/address accumulator ($32\times3\times
8 = 768$ words). Everything else --- hash intermediates, tournament condition
vectors --- comes from a small \emph{shared pool} of temporaries (sizes
$(16,4)$) that is reused group-to-group and wave-to-wave, rather than
allocated per group. Pool sizing is a direct scratch-vs-parallelism knob:
too small serializes on write-after-write hazards, too large overflows.

\paragraph{Constant coalescing.} The register-pressure analysis initially read
a peak of ${\sim}2300$ live words --- but ${\sim}2048$ of those were redundant
per-walker copies of a handful of \emph{compile-time constants} (gather-base
offsets, the reset position). Coalescing identical constants to one shared
broadcast word each --- what a real vectorized kernel does --- collapses the
genuine working set to ${\sim}780$ words, which is what makes the whole design
fit without spilling.

\paragraph{Shared tournament tables and in-\code{mem} priming.} The tree is
static, so each served level's value/difference tables are built into scratch
\emph{once} at setup and shared by all groups. For the constant-commuting
transform, the deep-level tree values are pre-xored with $C_6$ directly in
\code{mem} at setup (using the near-idle \code{store} engine during the ramp),
so no per-round scratch or work is spent maintaining the transformed domain.

\paragraph{Dead-register mining.} On the final round, 52 vectors are provably
dead (the position accumulators of groups that no longer gather, and the
node-value buffers of already-finished blocks). The drain-shortening
optimization reuses exactly those dead words to hold its extra tables, funding
a genuine improvement with \emph{zero} net scratch --- the only reason it fits
at 1{,}535/1{,}536.

\section{Lower bounds and the remaining gap}

The hash must be evaluated 4096 times; that work alone, on the 60-lane-op/cyc
\code{alu}${+}$\code{valu} pair, is a hard floor of ${\sim}\,780$ cycles.
Adding the index recurrence gives ${\sim}\,930$; adding the routing folds that
cannot fit on the 1-wide \code{flow} engine gives the realized \code{valu}
floor of ${\sim}\,1035$, and the kernel runs at 1{,}053 --- $1.7\%$ above it.
The gap from 930 to 1035 is an irreducible ``routing tax'' imposed by the
ISA: \code{flow} has only one slot per cycle, so the selects that would
otherwise offload from the saturated arithmetic engines serialize instead.
Going below 1{,}000 would require removing ${\sim}215$ \code{valu} slots of
real arithmetic --- i.e.\ a hash shorter than 11 operations (searched to
$2.4\times10^{12}$ candidates without a hit) or a structurally cheaper
routing scheme --- not better scheduling, which is already essentially
optimal.

\end{document}
